Um das ganze Eis zu schmelzen und das Schmelzwasser auf $\vartheta_\text{s} = \TsO$ zu erwärmen, braucht es mit $L_\text{f} = \LfEO$ (FoTaBe S.~190) und $c_\text{W} = \cWO$ (FoTaBe S.~191)
\begin{align}
 Q_\text{E} &= \QEF = \QE
\end{align}
Kondensiert der ganze Dampf, wird mit $L_\text{v} = \LvWO$ (FoTaBe S.~191) die Wärme
\begin{align}
 Q_\text{D} &= \QDF = \QD
\end{align}
frei, also viel mehr als nötig. Es kondensiert deshalb nur so viel Dampf, bis das Eis geschmolzen und das Wasser auf \TsO\ erwärmt ist:
\begin{align}
 m_\text{k} &= \mKoF = \frac{\QE}{\LvW} = \mKo \approx \result{\mKoP{-3}{2}}
\end{align}
Am Schluss liegen bei \TsO\ $m_\text{E}+m_\text{k} \approx \result{\mWsP{-3}{2}}$ Wasser und $m_\text{D}-m_\text{k} \approx \result{\mReP{-3}{2}}$ Wasserdampf vor: Das ganze Eis schmilzt, aber ein Teil des Wasserdampfs bleibt übrig.
